\section{Encoding, orbit structure, and cyclic accessibility}
\label{sec:rewrite-framework}

Let $\mathcal T=\{1,\ldots,n_T\}$ be a task universe and let
$\mathcal B_b\subseteq\mathcal T$, $1\leq b\leq n_B$, be the candidate
bases.  Equivalently, $A\in\{0,1\}^{n_B\times n_T}$ is their incidence
matrix, with $A_{bt}=1$ exactly when $t\in\mathcal B_b$.  Minimum Set Cover
asks for
\begin{equation}
 k^\star=\min\bigl\{|S|:S\subseteq[n_B],\quad
 \textstyle\bigcup_{b\in S}\mathcal B_b=\mathcal T\bigr\}.
 \label{eq:msc-definition}
\end{equation}
We fix $1\leq k\leq n_B$ address registers and first work in the physical
address space
\begin{equation}
 \mathcal H_{\rm addr}=(\mathbb C^{n_B})^{\otimes k}.
 \label{eq:addr-space-rewrite}
\end{equation}
The embedding into $2^{\lceil\log_2n_B\rceil}$ levels and its dynamically
disconnected padding block are implementation details.

Write $\ket{\boldsymbol b}=\ket{b_1}\otimes\cdots\otimes\ket{b_k}$ for a
computational-basis tuple.  Its numbers of uncovered tasks and repeated
base pairs are
\begin{align}
 U(\boldsymbol b)
 &=\sum_{t=1}^{n_T}\prod_{r=1}^k(1-A_{b_r,t}),\\
 D(\boldsymbol b)
 &=\sum_{1\leq r<s\leq k}\mathbf 1[b_r=b_s],
\end{align}
respectively, where $\mathbf 1[\cdot]$ is the indicator of the enclosed
condition.  Thus $U(\boldsymbol b)=D(\boldsymbol b)=0$ exactly when a
measurement returns an exact-$k$ cover with distinct bases.  When
$k=k^\star$, every such outcome solves MSC; for other $k$, the Hamiltonian
encodes the fixed-cardinality covering problem rather than minimum
cardinality itself.

Let $\ket e=\sum_{b=1}^{n_B}\ket b$ and let $I_B$ be the
identity on one address register.  The normalized complete-graph hopping
matrix and its $k$-register lift are
\begin{equation}
 W=\frac{\ket e\!\bra e-I_B}{n_B-1},
 \qquad
 H_{\rm walk}=\sum_{r=1}^k W^{(r)},
 \label{eq:walk-definition}
\end{equation}
where $W^{(r)}$ acts as $W$ on register $r$ and as the identity on all
others.  Define the diagonal penalty operators
\begin{equation}
 H_{\rm cov}=\sum_{\boldsymbol b}U(\boldsymbol b)
 \ket{\boldsymbol b}\!\bra{\boldsymbol b},
 \qquad
 H_{\rm excl}=\sum_{\boldsymbol b}D(\boldsymbol b)
 \ket{\boldsymbol b}\!\bra{\boldsymbol b}.
 \label{eq:penalty-definitions}
\end{equation}
Both sums run over all $n_B^k$ computational-basis tuples.
For positive penalty energies $\lambda$ and $\mu$, the problem Hamiltonian
is
\begin{equation}
 H_{\rm prob}=H_{\rm walk}+\lambda H_{\rm cov}+\mu H_{\rm excl},
 \label{eq:hprob-rewrite}
\end{equation}
so $\lambda U(\boldsymbol b)+\mu D(\boldsymbol b)$ is the total diagonal
penalty of $\ket{\boldsymbol b}$.  An explicit initial Hamiltonian is
\begin{equation}
 H_{\rm init}=\sum_{r=1}^k\bigl(I_B-\ket{u_B}\!\bra{u_B}\bigr)_r,
 \qquad
 \ket{u_B}=n_B^{-1/2}\sum_{b=1}^{n_B}\ket b,
 \label{eq:hinit-rewrite}
\end{equation}
with unique ground state $\ket u=\ket{u_B}^{\otimes k}$ and initial gap one.
The retained-walk linear interpolation is parameterized by $s\in[0,1]$ as
\begin{equation}
 H(s)=(1-s)H_{\rm init}+sH_{\rm prob}.
 \label{eq:cycle-linear-rewrite}
\end{equation}

\subsection{The faithfully represented group}

The colour-preserving incidence automorphism group is
\begin{equation}
 \begin{split}
 G_A=\{(\pi_B,\pi_T)\in {}&S_{n_B}\times S_{n_T}:\\
 &A_{\pi_B(b),\pi_T(t)}=A_{bt}\ \text{for all }b,t\},
 \end{split}
\end{equation}
where $S_d$ denotes the symmetric group on $d$ labels.
Only its base component acts on Eq.~\eqref{eq:addr-space-rewrite}.  We
therefore use the effective group
\begin{equation}
 G_B=\operatorname{pr}_B(G_A)\simeq G_A/\ker\rho,
 \qquad \rho(\pi_B,\pi_T)=\pi_B,
 \label{eq:faithful-base-group}
\end{equation}
where $\operatorname{pr}_B$ projects an automorphism onto its base
permutation and $\ker\rho$ contains the automorphisms acting trivially on
all bases.  Thus $G_B$ is the faithfully represented base group.  The
register and base actions commute, so the
represented symmetry is
\begin{equation}
 G=S_k\times G_B.
\end{equation}
Every term in Eqs.~\eqref{eq:hprob-rewrite} and \eqref{eq:hinit-rewrite}
commutes with $G$.  Consequently, any Hamiltonian $H$ commuting with $G$
has the block structure
\begin{align}
 \mathcal H_{\rm addr}
 &\simeq\bigoplus_{\alpha,\beta}
 V_\alpha^{S_k}\otimes V_\beta^{G_B}\otimes\mathcal M_{\alpha\beta},\\
 H&\simeq\bigoplus_{\alpha,\beta}
 I_{V_\alpha}\otimes I_{V_\beta}\otimes h_{\alpha\beta}.
 \label{eq:commutant-rewrite}
\end{align}
Here $\alpha$ and $\beta$ label irreducible representations of $S_k$ and
$G_B$; the spaces $V_\alpha^{S_k}$ and $V_\beta^{G_B}$ carry those
representations; $\mathcal M_{\alpha\beta}$ is their multiplicity space;
and $h_{\alpha\beta}$ is the reduced action on that space.  The symbols
$I_{V_\alpha}$ and
$I_{V_\beta}$ denote the corresponding identities.  This is the source of
Schur multiplicities; an equality of eigenvalues in
different blocks is a separate phenomenon.  The full representation and
the role of the kernel of $G_A$ are detailed in
Appendix~\ref{app:joint-representation}.

\subsection{Allowed space versus accessible space}

Let $U_\sigma$ permute the $k$ registers according to $\sigma\in S_k$, and
let $V_g$ apply the base permutation $g\in G_B$ in every register.  The
joint fixed projector and its range are
\begin{align}
 \Pi_{\rm sym}
 &=\left(\frac1{k!}\sum_{\sigma\in S_k}U_\sigma\right)
 \left(\frac1{|G_B|}\sum_{g\in G_B}V_g\right),\\
 \mathcal H_{\rm sym}&=\Ran\Pi_{\rm sym}.
 \label{eq:hsym-rewrite}
\end{align}
This is the symmetry-allowed invariant space, not automatically the space
filled by the protocol.  Let
$\mathfrak A=\operatorname{alg}^*(H_{\rm init},H_{\rm prob},I)$ be the
smallest unital algebra containing the endpoints and closed under adjoints;
here $I$ is the identity on $\mathcal H_{\rm addr}$.  The dynamically
accessible cyclic space is
\begin{multline}
 \mathcal K_u=\mathfrak A\ket u\\
 =\operatorname{span}\{X_\ell\cdots X_1\ket u:\
 X_i\in\{H_{\rm init},H_{\rm prob}\},\ \ell\geq0\}.
 \label{eq:ku-rewrite}
\end{multline}
It is the smallest common reducing subspace that contains the initial state,
and in general
\begin{equation}
 \mathcal K_u\subseteq\mathcal H_{\rm sym}
 \subseteq\mathcal H_{\rm addr}.
 \label{eq:nested-spaces-rewrite}
\end{equation}
The first inclusion can be strict; padding makes it strict automatically if
one defines the symmetry on the full qubit space.

Figure~\ref{fig:accessible-spectrum} summarizes the two distinctions used
throughout the paper.  Group invariance identifies an allowed space, while
the initial state and generator algebra select the smaller cyclic space.
An eigenvalue coincidence in another irrep may then close the global
multiplicity gap without closing the isolation gap seen by the protocol.

\begin{figure*}[t]
\centering
\begin{tikzpicture}[font=\small,>=Latex]
  \node[draw,rounded corners,minimum width=6.0cm,minimum height=3.1cm]
        (valid) {};
  \node[draw,rounded corners,fill=blue!5,minimum width=4.7cm,
        minimum height=2.1cm] (sym) at ([yshift=-0.25cm]valid.center) {};
  \node[draw,rounded corners,fill=blue!14,minimum width=2.9cm,
        minimum height=1.05cm] (cyc) at ([yshift=-0.25cm]sym.center) {};
  \node[anchor=north west] at ([xshift=0.12cm,yshift=-0.12cm]valid.north west)
        {\bfseries(a) $\mathcal H_{\rm addr}$};
  \node[anchor=north west] at ([xshift=0.12cm,yshift=-0.12cm]sym.north west)
        {$\mathcal H_{\rm sym}$};
  \node at (cyc.center) {$\mathcal K_u$};
  \node[align=center,below=0.12cm of valid]
        {symmetry allows $\mathcal H_{\rm sym}$;\\the protocol generates $\mathcal K_u$};

  \begin{scope}[xshift=6.0cm,yshift=-1.90cm]
    \draw[->] (0,0) -- (5.2,0) node[below] {$s$};
    \draw[->] (0,0) -- (0,3.2) node[left] {energy};
    \node[anchor=west] at (0.25,3.35) {\bfseries(b) Dark global crossing};
    \draw[very thick,blue] (0.25,0.55) .. controls (2.0,0.75) and
          (3.5,1.10) .. (4.85,1.25)
          node[right,blue,align=left] {tracked ground\\($\gamma$)};
    \draw[very thick,blue,dashed] (0.25,2.55) .. controls (2.0,2.35) and
          (3.5,2.25) .. (4.85,2.15)
          node[right,blue,align=left] {tracked excitation\\($\gamma$)};
    \draw[very thick,red!75!black] (0.25,2.35) .. controls (1.8,1.65) and
          (3.2,0.75) .. (4.85,0.25)
          node[right,red!75!black,align=left,yshift=0.12cm]
          {dark branch\\($\gamma'$)};
    \fill (3.05,1.00) circle (1.7pt);
    \draw[<->,blue] (2.25,0.86) -- (2.25,2.30)
          node[midway,right,align=left] {$\Delta_{\rm ad}>0$\\inside $\mathcal K_u$};
    \node[align=center,below=0.12cm] at (2.7,0)
          {different irreps do not hybridize};
  \end{scope}
\end{tikzpicture}
\caption{Protocol-dependent spectral relevance.  (a) The joint-fixed space
is only symmetry allowed; the cyclic space also depends on the initial state
and generated algebra.  (b) The labels $\gamma$ and $\gamma'$ denote
inequivalent symmetry sectors.  A branch in the latter can cross the followed
branch and close a global multiplicity gap while the tracked band remains
isolated within the cyclic space.  Solid blue is the tracked ground branch,
dashed blue its in-sector excitation, and red an inequivalent dark branch.}
\label{fig:accessible-spectrum}
\end{figure*}

\subsection{Orbit invariants and exact quotient}

For $\boldsymbol b=(b_1,\ldots,b_k)$, the joint action is
\begin{equation}
 \begin{aligned}
 (\sigma,g)\boldsymbol b
 &=\bigl(g(b_{\sigma^{-1}(1)}),\ldots,
 g(b_{\sigma^{-1}(k)})\bigr),\\
 &\hspace{4.5em}(\sigma,g)\in S_k\times G_B.
 \end{aligned}
 \label{eq:joint-tuple-action}
\end{equation}
The following elementary observation connects this representation directly
to the covering problem.

\begin{lemma}[Orbit invariance of the covering penalties]
\label{lem:orbit-invariance}
The functions $U$ and $D$ are constant on every orbit of the action in
Eq.~\eqref{eq:joint-tuple-action}.  Consequently,
\begin{equation}
 \Phi(\boldsymbol b)=\lambda U(\boldsymbol b)+\mu D(\boldsymbol b)
 \label{eq:orbit-potential}
\end{equation}
is an orbit function, the feasible set $\{U=D=0\}$ is a union of orbits,
and its projector $P$ commutes with $G$.
\end{lemma}

A register-permutation orbit is a size-$k$ multiset of bases.  A joint
$G$-orbit is therefore a $G_B$-equivalence class of such multisets: an
unlabelled candidate $k$-cover, which may still be infeasible or contain
repetitions.  The Cauchy--Frobenius lemma and its P\'olya--Redfield
cycle-index formulation~\cite{Burnside1911,Redfield1927,Polya1937} give
their number as
\begin{proposition}[Orbit classification and count]
\label{prop:orbit-count}
Let $c_\ell(g)$ be the number of length-$\ell$ cycles of $g\in G_B$ in its
base action.  Then
\begin{equation}
 \dim\mathcal H_{\rm sym}
 =\frac1{|G_B|}\sum_{g\in G_B}[z^k]
 \prod_{\ell\geq1}(1-z^\ell)^{-c_\ell(g)}.
 \label{eq:polya-rewrite}
\end{equation}
Here $[z^k]$ extracts the coefficient of $z^k$.  Equivalently, the
right-hand side is the cycle index of the base action evaluated on the
generating functions for multisets.
\end{proposition}

Equation~\eqref{eq:polya-rewrite} counts basis states in the group-fixed
space, not all vectors in the full quantum Hilbert space and not all states
that a particular protocol can prepare.  The latter question requires the
cyclic space below.

Orbit quotients of symmetry-confined quantum walks are established
constructions~\cite{KroviBrun2007}; we use the joint action here to retain
the exact Hamiltonian matrix elements.  For an orbit $\mathcal O$, define
\begin{equation}
 \ket{\mathcal O}=|\mathcal O|^{-1/2}
 \sum_{x\in\mathcal O}\ket x.
\end{equation}
These vectors are an orthonormal basis of $\mathcal H_{\rm sym}$.  If $H$
commutes with $G$, $x\in\mathcal O$, and
$H_{yx}:=\langle y|H|x\rangle$, then
\begin{equation}
 \langle\mathcal O'|H|\mathcal O\rangle
 =\sqrt{\frac{|\mathcal O|}{|\mathcal O'|}}
 \sum_{y\in\mathcal O'}H_{yx}.
 \label{eq:orbit-quotient-rewrite}
\end{equation}
The sum is independent of the representative $x$.  Lemma
\ref{lem:orbit-invariance} makes the penalty term diagonal in this basis:
\begin{equation}
 \langle\mathcal O'|V|\mathcal O\rangle
 =\delta_{\mathcal O\mathcal O'}\Phi(\mathcal O).
 \label{eq:orbit-potential-diagonal}
\end{equation}

Thus the quotient is a single-particle hopping problem.  Its sites are
unlabelled candidate covers, $\Phi$ is the on-site potential, and the
off-diagonal entries of Eq.~\eqref{eq:orbit-quotient-rewrite} are hopping
amplitudes with orbit-multiplicity factors.  The graph is connected because
one-register changes connect any two size-$k$ multisets.  Reachability of a
solution site is therefore not the obstruction: the protocol must
redistribute the state's amplitude into the fixed solution subspace
\begin{equation}
 \Ran P\cap\mathcal H_{\rm sym}
 =\operatorname{span}\{\ket{\mathcal O}:\Phi(\mathcal O)=0\}.
 \label{eq:solution-orbit-span}
\end{equation}
An individual computational solution need not be group fixed, but its
uniform orbit state always lies in Eq.~\eqref{eq:solution-orbit-span}.
Nontrivial combinations within the same solution orbit belong to other
irreps and can be inaccessible from $\ket u$.
Away from a hopping-cancellation point, orbit potentials are not spectral
gaps because eigenvectors are generally delocalized across several sites.
Moreover, nontrivial irreps have no site in this quotient; this is precisely
why inter-sector partners of dark crossings are absent from the dynamics
shown by the orbit graph.

\subsection{Cyclic accessibility in orbit coordinates}

Let
\begin{equation}
 M=\sum_{r=1}^k\ket e\!\bra e_{(r)},
 \qquad V=\lambda H_{\rm cov}+\mu H_{\rm excl}.
 \label{eq:transport-potential-generators}
\end{equation}
The endpoint algebra has only these two nontrivial generators.

\begin{lemma}[Two-generator reduction]
\label{lem:two-generator-reduction}
\begin{equation}
 \operatorname{alg}^*(H_{\rm init},H_{\rm prob},I)
 =\operatorname{alg}^*(M,V,I).
 \label{eq:two-generator-reduction}
\end{equation}
\end{lemma}

Let $\widehat M$ and $\widehat V$ be the orbit quotients of $M$ and $V$.
Write $\varphi_1<\cdots<\varphi_p$ for the distinct values of $\Phi$ on
the orbits and let $R_j$ project onto the corresponding potential shell.
The initial vector has strictly positive orbit coordinates
\begin{equation}
 w_{\mathcal O}=\langle\mathcal O|u\rangle
 =\sqrt{|\mathcal O|/n_B^k}>0.
 \label{eq:uniform-orbit-coordinates}
\end{equation}
Thus every orbit has nonzero amplitude initially, but this does not imply
that the individual basis vector $\ket{\mathcal O}$ can be prepared by the
protocol.  A strict cyclic inclusion excludes linear combinations---patterns
of relative amplitudes and phases across orbit sites---rather than deleting
particular sites from the graph.

\begin{theorem}[Shell-resolved cyclic closure]
\label{thm:shell-cyclic-closure}
In orbit coordinates, $\mathcal K_u$ is the smallest subspace containing
$w$ and invariant under $\widehat M$ and $\widehat V$.  It is obtained by
the stabilising iteration
\begin{align}
 \mathcal K^{(0)}&=\operatorname{span}\{R_jw:1\leq j\leq p\},\\
 \mathcal K^{(r+1)}&=\mathcal K^{(r)}+
 \sum_{j=1}^pR_j\widehat M\mathcal K^{(r)}.
 \label{eq:shell-cyclic-iteration}
\end{align}
Consequently,
\begin{equation}
 \dim\mathcal K_u\geq p,
 \label{eq:shell-lower-bound}
\end{equation}
and $\mathcal K_u=\mathcal H_{\rm sym}$ whenever $\Phi$ is injective on
orbits.  Conversely, a strict inclusion
$\mathcal K_u\subsetneq\mathcal H_{\rm sym}$ requires two distinct orbits
to share a potential value.

For the linear path in Eq.~\eqref{eq:cycle-linear-rewrite}, let
\begin{equation}
 c(s)=\frac{s}{n_B-1}-\frac{1-s}{n_B}.
 \label{eq:hopping-coefficient}
\end{equation}
Whenever $c(s)\leq0$, the symmetric and cyclic ground energies agree:
\begin{equation}
 E_0^{\mathcal K}(s)=E_0^{\rm sym}(s).
 \label{eq:zero-cyclic-cost-before-cancellation}
\end{equation}
\end{theorem}

The last statement follows physically from the sign-free, connected orbit
hopping problem.  Before cancellation, Perron--Frobenius gives a unique
strictly positive symmetric ground vector, which has nonzero overlap with
$w$ and hence belongs to the cyclic space.  At cancellation, a nonzero
projection of $w$ remains in the lowest potential shell.  Strict dimension
reduction therefore need not incur a ground-energy cost; the numerical
audit in Section~\ref{sec:numerical-reconstruction} finds precisely this
behaviour in four of the five strict inclusions.  Proofs of the orbit
statements and Theorem~\ref{thm:shell-cyclic-closure} are given in
Appendix~\ref{app:joint-representation}.
